\documentclass[../main.tex]{subfiles}
\begin{document}
\setlist[itemize]{topsep=2pt,itemsep=1pt,parsep=0pt,partopsep=0pt}
\setlist[itemize,2]{topsep=1pt,itemsep=0pt,parsep=0pt,partopsep=0pt}
\setlength{\parskip}{3pt}

\chapter{Week 2 — Research Design}
\label{ch:week2}

\textit{Snapshot and part summaries: see the overview on page~\pageref{ov:week2}.}

\section{Part 1 — Strategic Areas: Conceptual and Technical Design}
\label{sec:conceptual-technical}

\begin{itemize}
  \item \keyterm{Conceptual design}: what, why, and how much — four phases: objective, research framework, research question and conceptual model, operationalisation.
  \item \keyterm{Technical design}: how, where, and when — three phases: research strategy, research materials, research planning.
  \item Conceptual must precede technical; a method chosen before the objective is fixed serves no stated purpose.
\end{itemize}

\section{Part 2 — Conceptual Design: Objectives, Questions, and Operationalisation}
\label{sec:objectives-rqs}
\label{sec:operationalization}

\begin{itemize}
  \item \keyterm{Theory-oriented objectives}: theory development (build new theory) or theory validation (test a variable relationship).
  \item \keyterm{Practice-oriented objectives}: problem analysis, diagnosis, design, or change.
  \item \keyterm{Research objective} template: ``the objective is \textbf{(a)} the contribution, by \textbf{(b)} the method of providing it.''
  \item \keyterm{Criteria}: informative, useful, feasible, clear.
  \item \keyterm{Research questions}: open-ended (not yes/no), at least two sub-questions, yours to answer not the client's; each must be relevant, researchable, anchored, and precise.
  \item \keyterm{Research framework}: schematic sequencing the questions, naming required knowledge and methods (borrowed or developed).
  \item \keyterm{Operationalisation}: define key concepts exactly, state assumptions and what is neglected, name variables and how they are measured — the output that drives technical design (\autoref{sec:strategy-materials}).
\end{itemize}

\section{Part 3 — Technical Design: Strategy, Materials, and Planning}
\label{sec:strategy-materials}

\begin{itemize}
  \item \keyterm{Research strategy}: how the research will be done — judged on skill suitability, resource availability, achievable duration, and ethics.
  \item \keyterm{Research materials}: every tool required; secure access early, allow learning time, plan for failure points.
  \item \keyterm{Research planning}: time-based schedule; must allow flexibility. Includes milestones, deliverables, WBS, work-flow diagram, Gantt chart — developed further in Week~5 (\autoref{sec:project-mgmt}).
\end{itemize}

\section{Part 4 — Ethics and Research Conduct}
\label{sec:ethics}

\begin{itemize}
  \item Ensure reproducibility and traceability to original data; keep the raw data.
  \item Verify the origin of supplied data; give credit — unverified inputs and missing credit are the two most likely misconduct points.
  \item Follow the TU Delft Code of Conduct as the governing integrity standard.
  \item \keyterm{HREC}: any study involving human participants (questionnaires, interviews, test subjects) requires ethics board approval before it can proceed — plan for the approval timeline (\autoref{sec:data-sources}).
\end{itemize}

\section{Reading Material}
\label{sec:reading2}

\begin{partbox}
The Week~2 reading is the TU Delft Code of Conduct. It frames integrity as a shared institutional commitment built on six core values: diversity, integrity, respect, engagement, courage, and trust.
\end{partbox}

\begin{itemize}
  \item The six values underpin the Integrity Statement — basic principles every community member is expected to act upon.
  \item Integrity is a collective responsibility supported by institutional means, resources, and guidance.
\end{itemize}

\end{document}
